\documentclass[11pt]{article}
\usepackage[margin=1in]{geometry}
\usepackage{amsmath,amssymb,amsthm}
\usepackage[colorlinks=true,urlcolor=blue,linkcolor=blue,citecolor=blue]{hyperref}
\newtheorem{theorem}{Theorem}
\newtheorem{lemma}[theorem]{Lemma}
\newtheorem{corollary}[theorem]{Corollary}
\newcommand{\R}{\mathbb R}
\newcommand{\norm}[1]{\left\lVert#1\right\rVert}
\newcommand{\K}{\mathcal K}
\newcommand{\range}{\operatorname{range}}
\title{Direct cluster-robust bounds for randomized\\small-block Krylov spaces}
\author{Working manuscript: proof and priority review in progress}
\date{November 14, 2026}
\begin{document}
\maketitle
\begin{abstract}
We bound the actual fundamental matrix interpolation polynomials for
independent Gaussian similarities of real diagonal blocks. For every block
size $b$, number of blocks $d$, and fixed admissible spectra in an interval
of length $L$, their simultaneous growth is at most
\[
 (1000d^8b^5\beta^{-4}L/\Delta)^{d-1}
\]
with probability at least $1-\beta$.
Here $\Delta$ is the minimum cross-block gap; no within-block gap appears.
A terminal block-Lanczos polynomial has a resolvent inverse bounded by the
norm of the original Gaussian start. Conditioning on the other blocks then
restores independent full-covariance Gaussian rows at a fresh block.
Fractional moments along elimination paths keep the logarithm of the
prefactor linear in $d$, up to logarithms. This proves a direct
cluster-robust small-block Krylov angle bound and its scalar-filtered
consequence. For positive target spectra we additionally prove uniform
growth on the lower tail in terms of relative cross-block gaps, removing
condition-number dependence. Analysis-only padding extends this to a
sliding $b$th-order relative gap, at at most twice the startup depth.
It does not prove the previously conjectured intermediate
factor bounds, which have a known counterexample. No new matrix-product
quadratic-time algorithm or finite-precision theorem is claimed.
\end{abstract}

\section{Statement and scope}
For a matrix polynomial $P(z)=\sum_{j=0}^{s}z^j C_j$, always use
\emph{left evaluation with right coefficients}:
\[
 P(B)=\sum_{j=0}^{s}B^j C_j.
\]
Scalar-variable products cannot generally be evaluated by substituting a
matrix into each factor separately.

Fix integers $b\geq1$, $d\geq2$. Let $\Lambda_1,\ldots,\Lambda_d$
be fixed real diagonal $b$-by-$b$ matrices with pairwise disjoint spectra.
Draw independent standard Gaussian $\Omega_i\in\R^{b\times b}$ and put
$B_i=\Omega_i^{-1}\Lambda_i\Omega_i$. All inverses in the statements
exist almost surely. Let $[a,c]$ contain the spectra, $L=c-a>0$, and
\[
 \Delta=\min_{i\ne j}\min_{\lambda\in\operatorname{spec}\Lambda_i,
                 \eta\in\operatorname{spec}\Lambda_j}|\lambda-\eta|>0.
\]
Write $F_i$ for the degree-at-most-$d-1$ fundamental polynomials satisfying
$F_i(B_j)=\delta_{ij}I_b$.

\begin{theorem}[Direct interpolation growth]\label{thm:main}
For every $0<\beta<1$ and every fixed admissible collection of spectra,
\[
 \Pr\left\{\max_i\sup_{z\in[a,c]}\norm{F_i(z)}_2
 \leq\left(\frac{1000d^8b^5}{\beta^4}\frac L\Delta\right)^{d-1}
 \right\}\geq1-\beta.
\]
The constant is independent of all eigenvalues and all within-block gaps,
including zero within-block gaps. In particular, the logarithm of the
spectrum-independent prefactor is $O(d\log(db/\beta))$.
\end{theorem}

The block partition is part of the hypothesis. Positive cross-block gap
for a partition into equal-size blocks is not equivalent to saying only
that each eigenvalue has multiplicity at most $b$.
Probability is separate for every fixed input. The theorem does not say
that one random draw works for every spectrum simultaneously.

Shao's structural analysis \cite{shao} formulates the corresponding Krylov
angle question using intermediate nonnormal factors. Their proposed
spectrum-independent bounds were subsequently refuted \cite{stepaniants}.
Theorem~\ref{thm:main} controls the actual interpolants instead. Its
possible novelty and precise relationship to other cluster-robustness
theorems are under independent review; none of the classical ingredients
below is individually claimed new.

\section{Two elementary preliminaries}
\begin{lemma}[Independent Gaussian rows]\label{lem:rows}
Let $N\in\R^{b\times b}$ have independent centered Gaussian rows whose
covariance matrices are each at least $h^2I_b$, where $h>0$. Then
\[
 \Pr\{\sigma_{\min}(N)\leq ht\}\leq b^{3/2}t\qquad(t>0).
\]
The assertion also holds conditionally when the covariance lower bound is
fixed by the conditioning.
\end{lemma}
\begin{proof}
Conditioning on the other rows, a row's distance $D_i$ to their span has
small-ball probability at most $\sqrt{2/\pi}\,u/h$ at radius $u$.
The other rows have rank $b-1$ almost surely. For invertible $N$,
$\norm{N^{-1}}_F^2=\sum_iD_i^{-2}$. If all $D_i>\sqrt b\,ht$, this
sum is less than $(ht)^{-2}$. A union bound proves the claim. The same
argument proves nonsingularity almost surely and covers $b=1$.
\end{proof}

\begin{lemma}[Generic interpolation]\label{lem:generic}
For any nonempty subset $S$ of the blocks, let
$A_S=\operatorname{diag}(\Lambda_i:i\in S)$,
$\Gamma_S=\operatorname{col}(\Omega_i:i\in S)$, and $s=|S|$.
Then
\[
 [\Gamma_S,A_S\Gamma_S,\ldots,A_S^{s-1}\Gamma_S]
\]
is nonsingular almost surely. Consequently the degree-at-most-$s-1$
interpolation map on the matrices $(B_i:i\in S)$ is bijective.
These assertions hold simultaneously for all subsets.
\end{lemma}
\begin{proof}
The determinant is a polynomial in the Gaussian entries. At every
$\Omega_i=I_b$, permutation of rows and columns splits it into $b$
scalar Vandermonde determinants, each using one eigenvalue from each
block. All are nonzero by cross-block disjointness. Thus its zero set
has Gaussian probability zero. Multiplication of block row $i$ by
$\Omega_i^{-1}$ identifies this matrix with the left-evaluation map.
There are only finitely many subsets.
\end{proof}

\section{A terminal polynomial with a controlled resolvent inverse}
We normalize all spectra to $[0,1]$ until the end of the proof.
Consider a subset of $q\geq1$ blocks, and abbreviate $A=A_S$ and
$\Gamma=\Gamma_S$. Choose a block orthonormal Krylov basis
$V_0,\ldots,V_{q-1}$ by a fixed measurable Gram--Schmidt/QR rule applied
to $[\Gamma,A\Gamma,\ldots,A^{q-1}\Gamma]$. This rule uses only the
blocks in $S$. Lemma~\ref{lem:generic} rules out breakdown.
There are right-coefficient polynomials $P_j$ of degree $j$ such that
\[
 V_j=\sum_{t=0}^j A^t\Gamma C_{t,j},\qquad
 P_j(z)=\sum_{t=0}^jz^t C_{t,j},
\]
with nonsingular leading coefficient $C_{j,j}$.
By symmetry of $A$, the terminal block-Lanczos recurrence is
\[
 A V_{q-1}=V_{q-2}E_{q-1}^{\mathsf T}+V_{q-1}D_{q-1}.
\]
The first term is absent for $q=1$. Here $\norm{D_{q-1}}_2\leq1$,
$\norm{E_{q-1}}_2\leq1$, since $\norm A_2\leq1$.
Define
\begin{equation}\label{eq:R}
 R(z)=zP_{q-1}(z)-P_{q-2}(z)E_{q-1}^{\mathsf T}
                         -P_{q-1}(z)D_{q-1}.
\end{equation}

\begin{lemma}[Terminal resolvent identity]\label{lem:resolvent}
The polynomial $R$ annihilates each $B_i$, $i\in S$, under left
evaluation. For $z\notin\operatorname{spec}(A)$,
\[
 W(z):=V_{q-1}^{\mathsf T}(zI-A)^{-1}\Gamma,
 \qquad W(z)R(z)=I_b.
\]
In particular $R(z)$ is invertible there and
\[
 \norm{R(z)^{-1}}_2\leq
 \frac{\norm\Gamma_2}{\operatorname{dist}(z,\operatorname{spec}A)}.
\]
If $G_S=\max_{i\in S}\sup_{z\in[0,1]}\norm{F_i^S(z)}_2$ denotes
the actual fundamental-polynomial growth on $S$, then
\[
 \sup_{z\in[0,1]}\norm{R(z)}_2
       \leq3qG_S\max_{i\in S}\norm{\Omega_i^{-1}}_2.
\]
\end{lemma}
\begin{proof}
The terminal recurrence says $\sum_{\ell=0}^qA^\ell\Gamma C_\ell=0$
for the coefficients of $R$. Its block rows imply $R(B_i)=0$.
For every integer $\ell\geq1$,
\[
 (zI-A)^{-1}(z^\ell\Gamma-A^\ell\Gamma)
       =\sum_{t=0}^{\ell-1}z^{\ell-1-t}A^t\Gamma.
\]
Multiply by $C_\ell$, sum, and use annihilation. Multiplication on the
left by $V_{q-1}^{\mathsf T}$ kills every term with $t<q-1$.
Only $\ell=q,t=q-1$ remains, and it equals $I_b$ because
\[
 V_{q-1}^{\mathsf T}A^{q-1}\Gamma C_q=I_b.
\]
Indeed $C_q=C_{q-1,q-1}$ and $V_{q-1}$ is its leading Krylov term plus
terms orthogonal to $V_{q-1}$. This proves $WR=I$. A square matrix with
a right inverse is invertible. The resolvent bound follows from symmetry
of $A$ and $\norm{V_{q-1}}_2=1$.

For the last claim, uniqueness of interpolation gives
\[
 P_j(z)=\sum_{i\in S}F_i^S(z)\Omega_i^{-1}(V_j)_i.
\]
Each block row $(V_j)_i$ has norm at most one. Thus
$\sup\norm{P_j}\leq qG_S\max_i\norm{\Omega_i^{-1}}$.
Apply \eqref{eq:R}, $|z|\leq1$, and the recurrence-coefficient bounds.
This also proves the $q=1$ case with its absent previous block.
\end{proof}

\section{A fresh Gaussian block and the deterministic recursion}
Let $T$ be a subset of at least two blocks and choose $k\in T$.
Apply Lemma~\ref{lem:resolvent} on $S=T\setminus\{k\}$, independently
of $\Omega_k$, and define
\[
 N_{T,k}=\Omega_k R(B_k),\qquad
 \delta_{T,k}=\min_{i\in S}\operatorname{dist}
                 (\operatorname{spec}\Lambda_i,\operatorname{spec}\Lambda_k).
\]
Conditional on $(\Omega_i:i\in S)$, row $\ell$ of $N_{T,k}$ is
$\omega_{k,\ell}^{\mathsf T}R(\lambda_{k,\ell})$.
The rows are independent centered Gaussians. Their covariance matrices
are each at least $(\delta_{T,k}/\norm{\Gamma_S}_2)^2I_b$.
Consequently $N_{T,k}$ is nonsingular almost surely and
\begin{equation}\label{eq:fresh}
 F_k^T(z)=R(z)N_{T,k}^{-1}\Omega_k.
\end{equation}
Its values are zero on $S$ and $I_b$ at $B_k$, so this is the actual
fundamental polynomial, not a different interpolation convention.

Define the random variable
\[
 Z_{T,k}=\frac{\delta_{T,k}\norm{N_{T,k}^{-1}}_2}
                      {\norm{\Gamma_S}_2}.
\]
Lemma~\ref{lem:rows} proves the conditional tail and fractional moment
\begin{equation}\label{eq:moment}
 \Pr\{Z_{T,k}>z\mid(\Omega_i:i\in S)\}
       \leq\min\{1,b^{3/2}/z\},\qquad
 \mathbb E[Z_{T,k}^{1/2}\mid(\Omega_i:i\in S)]\leq2b^{3/4}.
\end{equation}
The second bound is the integral of the first, split at $z=b^{3/2}$.
These are unconditional-on-success estimates: we do not condition them
on other roots' good events or on Gaussian norm bounds.

Let $\mathcal E$ be the simultaneous event
\[
 \norm{\Omega_i}_2\leq L_0=2b\sqrt{d/\beta},\qquad
 \norm{\Omega_i^{-1}}_2\leq U_0=4db^{3/2}/\beta
             \quad(1\leq i\leq d).
\]
Frobenius Markov bounds and Lemma~\ref{lem:rows} give
$\Pr(\mathcal E^c)\leq\beta/2$.
On $\mathcal E$, $\norm{\Gamma_S}_2\leq\sqrt qL_0$, $q=|T|-1$.
Equations \eqref{eq:fresh} and Lemma~\ref{lem:resolvent} therefore give
\begin{equation}\label{eq:recursion}
 G_T\leq\frac{3q^{3/2}U_0L_0^2}{\Delta}
                    \max_{k\in T}\{Z_{T,k}G_{T\setminus\{k\}}\}.
\end{equation}
For a singleton, $G_T=1$.

\section{Fractional moments along elimination paths}
For a permutation $\pi$ of the $d$ indices, put
$T_j=\{\pi_j,\ldots,\pi_d\}$ and
\[
 Y_\pi=\prod_{j=1}^{d-1}Z_{T_j,\pi_j}.
\]
All factors after the first are measurable in the Gaussian matrices of
$T_2$. Thus the tower property and \eqref{eq:moment}, successively
integrating out $\Omega_{\pi_1},\Omega_{\pi_2},\ldots$, give
\[
 \mathbb E Y_\pi^{1/2}\leq(2b^{3/4})^{d-1}.
\]
No independence between the factors is asserted or required. Markov's
inequality and a union bound over the $d!$ permutations show, with failure
at most $\beta/2$, that
\[
 \max_\pi Y_\pi\leq
 \left(\frac{2d!(2b^{3/4})^{d-1}}\beta\right)^2.
\]
Expand the deterministic recursion \eqref{eq:recursion} on $\mathcal E$.
Every elimination path has the same product of the factors $q^{3/2}$.
Combining the two events, with probability at least $1-\beta$,
\begin{align*}
 G_{\{1,\ldots,d\}}
 &\leq4\beta^{-2}(d!)^2((d-1)!)^{3/2}
       \left(\frac{192d^2b^5}{\beta^2\Delta}\right)^{d-1}\\
 &\leq\left(\frac{768d^8b^5}{\beta^4\Delta}\right)^{d-1}.
\end{align*}
For the second inequality one may use
$(d!)^2\leq d^{4(d-1)}$, $((d-1)!)^{3/2}\leq d^{3(d-1)/2}$,
and $4\beta^{-2}\leq(4\beta^{-2})^{d-1}$ for $d\geq2$.
The harmless round-up to $1000$ gives Theorem~\ref{thm:main} on $[0,1]$.
Affine scaling from $[a,c]$ replaces $\Delta$ by $\Delta/L$ and preserves
the left interpolation conditions. This completes its proof.

\section{Krylov angles and standard scalar acceleration}
Let $A\in\R^{n\times n}$ be fixed symmetric. Let $Q$ and $Q_\perp$ be
orthonormal eigenvector matrices whose columns together form an orthogonal
matrix. The target invariant subspace $\range Q$ has dimension $bd\leq n$,
and its eigenvalues admit the block partition
above. Let $L=\lambda_{\max}(A)-\lambda_{\min}(A)>0$ and
$g=\Delta/L$. Draw a standard Gaussian $\Omega\in\R^{n\times b}$.
Write $\Omega_H=Q^{\mathsf T}\Omega=\operatorname{col}(\Omega_i)$ and
$\Omega_T=Q_\perp^{\mathsf T}\Omega$. These matrices have independent
standard Gaussian entries in this fixed eigenbasis. Define
\[
 \K_s(A,\Omega)=\range[\Omega,A\Omega,\ldots,A^{s-1}\Omega].
\]

\begin{corollary}[Cluster-robust angle bound]\label{cor:angle}
For $n>bd$ and $0<\beta<1$, with probability at least $1-\beta$,
\[
 \tan\angle(\range Q,\K_d(A,\Omega))
 \leq\frac{8db^{3/2}}\beta
       \sqrt{\frac{2db(n-bd)}\beta}
       \left(\frac{16000d^8b^5}{\beta^4g}\right)^{d-1}.
\]
The angle is the largest principal angle. Within-block gaps do not enter.
If $n=bd$, the angle is zero almost surely.
\end{corollary}
\begin{proof}
Let $K_H=[\Omega_H,\Lambda_H\Omega_H,\ldots,
\Lambda_H^{d-1}\Omega_H]$, and define $K_T$ analogously.
For a partitioned unit head vector $y=\operatorname{col}(y_i)$,
\[
 p_y(z)=\sum_iF_i(z)\Omega_i^{-1}y_i
\]
is its unique interpolation polynomial. Its coefficients are $K_H^{-1}y$.
For every tail row at eigenvalue $\lambda_\ell$,
\[
 |\omega_\ell^{\mathsf T}p_y(\lambda_\ell)|
 \leq\norm{\omega_\ell}_2\sqrt d\,G_H
                     \max_i\norm{\Omega_i^{-1}}_2.
\]
Thus $\norm{K_TK_H^{-1}}_2\leq
\sqrt d\norm{\Omega_T}_F G_H\max_i\norm{\Omega_i^{-1}}_2$.
The graph-subspace formula identifies this norm with the tangent.
Use the combined event constructed in the proof of the head theorem
at failure $\beta/2$, including its Gaussian norm and inverse bounds,
giving inverse bound $8db^{3/2}/\beta$. Independently a tail Frobenius
Markov bound gives $\norm{\Omega_T}_F\leq
\sqrt{2b(n-bd)/\beta}$ except probability $\beta/2$.
\end{proof}

For any scalar polynomial $f$ of degree $t$ that is nonzero on the
target spectrum, the same graph argument and
$f(A)\K_d(A,\Omega)\subseteq\K_{d+t}(A,\Omega)$ give
\begin{equation}\label{eq:filtered}
 \tan\angle(\range Q,\K_{d+t}(A,\Omega))
 \leq\frac{\max_{\lambda\in\operatorname{spec}\Lambda_T}|f(\lambda)|}
              {\min_{\lambda\in\operatorname{spec}\Lambda_H}|f(\lambda)|}
             \norm{K_TK_H^{-1}}_2.
\end{equation}
This uses ordinary block Krylov iteration, without constructing the
interpolation polynomials in an implementation. Standard Chebyshev filters
therefore convert the initialization estimate into an accuracy bound.
For example, for the top $bd$ eigenvalues, write $\mu=\lambda_{bd+1}$,
$\ell=\lambda_n<\mu$, and
$\gamma=(\lambda_{bd}-\mu)/(\mu-\ell)>0$. The degree-$t$ Chebyshev
polynomial on $[\ell,\mu]$ has ratio in \eqref{eq:filtered} at most
$2\exp[-t\operatorname{arcosh}(1+2\gamma)]$.
If $C_\beta$ denotes the displayed right side in
Corollary~\ref{cor:angle}, a sufficient additional depth is
$t\geq\log(2C_\beta/\varepsilon)/\operatorname{arcosh}(1+2\gamma)$.
Its logarithmic numerator includes $\log(n/(\beta\varepsilon))$,
$d\log(1/g)$, and $d\log(db/\beta)$, with no within-block separation
cost. If the tail is a single eigenvalue, a
degree-one scalar factor annihilates it instead.

\section{Relative tail growth without condition-number dependence}
The global interval in Theorem~\ref{thm:main} is unnecessary when only
lower-tail evaluations are needed. Assume all target nodes are positive,
and put
\[
 h=\min_i\lambda_{\min}(\Lambda_i)>0,\qquad
 g_{\rm rel}=\min_{i\ne j}\min_{\lambda\in\operatorname{spec}\Lambda_i,
                   \eta\in\operatorname{spec}\Lambda_j}
              \frac{|\lambda-\eta|}{\max(\lambda,\eta)}>0.
\]
No ratio between the largest and smallest target eigenvalues will enter.

\begin{lemma}[Normalized resolvent identity]\label{lem:relative}
Use the terminal polynomial of Lemma~\ref{lem:resolvent} on any subset
$S$ of $q$ blocks. For $z\notin\operatorname{spec}(A_S)$ and
$\eta\notin\operatorname{spec}(A_S)$,
\[
 R(z)R(\eta)^{-1}
 =\sum_{i\in S}F_i^S(z)\Omega_i^{-1}
       (zI-\Lambda_i)(\eta I-\Lambda_i)^{-1}\Omega_i.
\]
In particular, if $0\leq z<h$ and $\eta$ belongs to a fresh target
block, then
\[
 \norm{R(z)R(\eta)^{-1}}_2
       \leq\frac{qG_S(z)}{g_{\rm rel}}
                              \max_{i\in S}\kappa(\Omega_i),
 \quad G_S(z)=\max_{i\in S}\norm{F_i^S(z)}_2.
\]
\end{lemma}
\begin{proof}
Let $A=A_S$, $\Gamma=\Gamma_S$, and
$\mathcal L_z=R(z)V_{q-1}^{\mathsf T}(zI-A)^{-1}$.
Polynomial division and $V_{q-1}\perp\K_{q-1}(A,\Gamma)$ give
\[
 \mathcal L_z A^t\Gamma=z^t I_b\qquad(0\leq t<q).
\]
Thus $\mathcal L_z$ is the exact evaluation-at-$z$ functional on the
invertible raw Krylov matrix. Uniqueness of interpolation identifies its
block row $i$ with $F_i^S(z)\Omega_i^{-1}$. Therefore
\[
 R(z)R(\eta)^{-1}
   =\mathcal L_z(zI-A)(\eta I-A)^{-1}\Gamma,
\]
which is the displayed formula. The unscaled symmetric operator may be
used here; alternatively affine rescaling of the preceding lemma cancels
in the normalized identity. For a scalar target node $\lambda$,
$0\leq z<h$ implies $|z-\lambda|\leq\lambda$, whereas
$|\eta-\lambda|\geq g_{\rm rel}\max(\eta,\lambda)$.
Diagonal similarity proves the norm bound.
\end{proof}

\begin{theorem}[Uniform relative-gap bound on the lower tail]
\label{thm:relative}
Under these positive-node assumptions, for every $0<\beta<1$,
\[
 \Pr\left\{\max_i\sup_{0\leq z\leq h}\norm{F_i(z)}_2
 \leq\left(\frac{20000d^{12}b^5}{\beta^4g_{\rm rel}}\right)^{d-1}
 \right\}\geq1-\beta.
\]
The good event can simultaneously include the Gaussian norm and inverse
bounds $L_0,U_0$ used in the proof of Theorem~\ref{thm:main}.
\end{theorem}
\begin{proof}
First fix a finite set $\mathcal Z\subset[0,h)$ of $N$ evaluation points.
For each subset $T$, root $k$, and $z\in\mathcal Z$, construct the
terminal $R$ on $S=T\setminus\{k\}$ as before, and now normalize by
\[
 \widetilde N_{T,k,z}=\Omega_k R(B_k)R(z)^{-1}.
\]
The fresh Gaussian rows remain conditionally independent. By
Lemma~\ref{lem:relative}, their covariance matrices have the common
lower bound
\[
 \left(\frac{g_{\rm rel}}{qG_S(z)\max_{i\in S}\kappa(\Omega_i)}\right)^2 I.
\]
Define $Z_{T,k,z}$ to be the norm of $\widetilde N_{T,k,z}^{-1}$
times the scalar inside these parentheses. It obeys the same conditional
tail and square-root moment bounds as \eqref{eq:moment}.
Moreover $F_k^T(z)=\widetilde N_{T,k,z}^{-1}\Omega_k$.
On the Gaussian norm event the deterministic recursion becomes
\[
 G_T(z)\leq\frac{qU_0L_0^2}{g_{\rm rel}}
             \max_{k\in T}\{Z_{T,k,z}G_{T\setminus\{k\}}(z)\}.
\]
Here $G_S(z)>0$ because the fundamental polynomials sum to $I$.

Apply the same conditional fractional-moment argument along every
elimination path, at each of the $N$ fixed points. Unioning only at the
\emph{final} Markov step over $Nd!$ possibilities gives, jointly with
the Gaussian norm event and with total failure at most $\beta$,
\begin{equation}\label{eq:relative-grid}
 \max_{z\in\mathcal Z}G_{\{1,\ldots,d\}}(z)
 \leq4N^2\beta^{-2}(d!)^2(d-1)!
       \left(\frac{64d^2b^5}{\beta^2g_{\rm rel}}\right)^{d-1}.
\end{equation}
There is no division of the failure budget by $N$ at every depth.

Finally put $p=d-1$, $N=\lceil4p^2\rceil+1$, and
$z_j=jh/N$ for $0\leq j<N$. All points are strictly below $h$, so no
normalization at a root was used. Scalar Markov's polynomial inequality,
applied to $u^{\mathsf T}F_i(z)v$ for arbitrary unit $u,v$, gives
\[
 \sup_{[0,h]}\norm{F_i'}_2
       \leq\frac{2p^2}{h}\sup_{[0,h]}\norm{F_i}_2.
\]
Every point of $[0,h]$ is within $h/N$ of a grid point. Thus
$\sup\norm{F_i}\leq2\max_j\norm{F_i(z_j)}$.
Multiply \eqref{eq:relative-grid} by two, use $N\leq5d^2$, and bound
the factorials as before. The resulting constant is at most
$(12800d^{11}b^5\beta^{-4}/g_{\rm rel})^{d-1}$, so the stated
round-up is valid.
\end{proof}

For a PSD operator whose top $bd$ eigenvalues have this partition,
all remaining eigenvalues lie in $[0,h]$. Repeating the graph proof of
Corollary~\ref{cor:angle} with Theorem~\ref{thm:relative} yields
\begin{equation}\label{eq:relative-angle}
 \tan\angle(Q,\K_d(A,\Omega))
 \leq\frac{8db^{3/2}}\beta
       \sqrt{\frac{2db(n-bd)}\beta}
       \left(\frac{320000d^{12}b^5}{\beta^4g_{\rm rel}}\right)^{d-1}
\end{equation}
with probability at least $1-\beta$. In particular the logarithm of
this bound is
\[
 O\left(d\log\frac{db}{\beta g_{\rm rel}}+\log\frac n\beta\right),
\]
with no target condition number.

\section{Sliding \texorpdfstring{$b$}{b}th gaps, packing, and analysis-only padding}
Suppose the top $k>b$ eigenvalues of a PSD operator are positive and set
\[
 \gamma_b=\min_{1\leq i\leq k-b}
                  \frac{\lambda_i-\lambda_{i+b}}{\lambda_i}>0.
\]
They need not admit a partition into exactly $k/b$ equal-size blocks.
The following construction avoids that limitation.

Connect adjacent ordered target eigenvalues whenever
$(\lambda_i-\lambda_{i+1})/\lambda_i<\gamma_b/b$.
Each connected component contains at most $b$ eigenvalues. Otherwise
a chain of $b$ such steps would give
\[
 \lambda_{i+b}>(1-\gamma_b/b)^b\lambda_i
                       \geq(1-\gamma_b)\lambda_i,
\]
contrary to the definition of $\gamma_b$.
Pack whole consecutive components into blocks of size at most $b$ by
the next-fit rule. Consecutive packed block sizes sum to more than $b$,
so the number $D$ of packed blocks is at most
$2\lceil k/b\rceil-1$. Every cross-block relative gap is at least
$\gamma_b/b$: the smallest relative separation is attained at a
boundary, and every such boundary is an unconnected adjacent pair.
Pad each block to size $b$ by repeating a node already in that block.
This preserves cross-block disjointness and all positive-node bounds.

Padding is only an analysis device. Add the corresponding diagonal
coordinates to $A$ and independent Gaussian rows to $\Omega$.
The enlarged head has dimension $Db$, but the tail is unchanged.
For any original head vector $y$, prescribe augmented head values
$(y,0_{\rm dummy})$ and use the invertible augmented head Krylov matrix.
Projection onto original coordinates produces a vector of the original
$\K_D(A,\Omega)$ with original head $y$ and tail norm bounded by the
augmented graph bound. Hence its largest principal angle is no larger.
The resulting property depends only on the original Gaussian start;
its marginal probability is therefore at least the joint success
probability. No dummy samples or target partition are supplied to the
actual algorithm.

Consequently \eqref{eq:relative-angle}, with $d=D$,
$g_{\rm rel}\geq\gamma_b/b$, and the enlarged dimension in its harmless
prefactor, proves a startup angle whose logarithm is
\begin{equation}\label{eq:packed-log}
 O\left(D\log\frac{Db}{\beta\gamma_b}+\log\frac n\beta\right),
 \qquad D\leq2\lceil k/b\rceil-1.
\end{equation}
The case $k\leq b$ is the usual Gaussian large-block initialization.

For low-rank approximation of a rectangular matrix $M$ and
$0<\varepsilon<1$, apply this to
$A=MM^{\mathsf T}$, $\lambda_i=\sigma_i(M)^2$. The established
good-start block-Krylov theorem \cite{mmm,chen} converts an initial
angle $T$ into rank-$k$ relative spectral and Frobenius approximation
after $O(\varepsilon^{-1/2}\log(n(1+T^2)/\varepsilon))$ further
iterations. To see the exact good-start property for the possibly wide
startup space, the graph construction supplies $Y=[I_k;X]$ in that
space, with $\norm X_2\leq T$. Its orthonormalization
$Y(I_k+X^{\mathsf T}X)^{-1/2}$ has inverse head-compression norm squared
$1+\norm X_2^2\leq1+T^2$. This is the existence-based
$(k,1+T^2)$-good property used in \cite[Definition 3.1]{chen}.
The imported theorem applies to an orthonormal basis of the entire
initial space; the number of its columns need not equal $k$.
There is the exact span identity
\[
 \K_t(A,\operatorname{orth}\K_D(A,\Omega))
                         =\K_{t+D-1}(A,\Omega).
\]
Thus ordinary randomized block Krylov has the sufficient depth
\begin{equation}\label{eq:lra-depth}
 O\left(\frac{D\log(Db/(\beta\gamma_b))
                         +\log(n/(\beta\varepsilon))}{\sqrt\varepsilon}\right).
\end{equation}
This uses the sliding relative $b$th gap and no singular-value condition
number. Matrix-vector count is $b$ times the depth. The usual
gap-dependent good-start theorem similarly replaces $\sqrt\varepsilon$
by $\sqrt{(\lambda_k-\lambda_\ell)/\lambda_k}$, with initialization
through $\ell>k$ positive target eigenvalues and their sliding gap.
When the optimal rank-$k$ tail is zero, multiplying the padded startup
space by $A$ annihilates the zero tail and gives the whole range after
one additional iteration, so this branch is exact.

These are consequences for the existing orthogonalized algorithm, not a
new expensive interpolation implementation. They address the
condition-number-free convergence question in \cite[Section 5]{chen}
up to additional logarithmic factors and the constant-factor startup
padding. They do not prove its literal unnormalized monomial
smallest-singular-value conjecture or assert the exact same quantitative
formula without those extra logarithms.

The zero-tail branch above uses this section's positive-top-$k$
hypothesis, so zero tail means that the positive rank is exactly $k$.
It is not a claim that an arbitrarily large repeated zero eigenspace can
be recovered from a smaller block.

\paragraph{Ordinary arithmetic realization.}
These bounds concern the existing orthogonalized algorithm. With
$w=bq$ raw basis columns, standard dense operations cost
$O(npw+(n+p)w^2+w^3)$ for an explicit $n$-by-$p$ matrix, including
orthogonalization and compressed SVD. For a factored matrix
$M=A_1B_1$, of inner dimension $m\leq\min(n,p)$ and input size
$I=m(n+p)$, factor actions replace $npw$ by $Iw$.
When $w\leq m$, all displayed basis and core terms are at most $O(Iw)$.
Expansion of the final rank-$k$ factors, if requested, additionally costs
$O(npk)$. Larger bases may use the usual exact QR/core-SVD fallback.
No asymptotically fast multiplication or new numerical primitive is
needed. The generic $\varepsilon^{-1/2}$ accuracy dependence remains;
the new assertion concerns clustering and condition-number dependence.

\section{Why raw monomial conditioning is a different question}
For completeness, a uniform lower bound for the smallest singular value
of a raw monomial Krylov matrix cannot discard the condition number
while keeping the stated exponential dependence on depth. Take $b=1$,
a Gaussian start $g$, and
\[
 A=\operatorname{diag}(1,2^{-1},\ldots,2^{-(k-1)}).
\]
The polynomial
\[
 p(x)=\prod_{j=0}^{k-2}(x-2^{-j})
\]
has coefficient norm at least one. For
$K=[g,Ag,\ldots,A^{k-1}g]$ its image has only a last coordinate, so
\[
 \sigma_{\min}(K)\leq|g_k|\,2^{-(k-1)(k-2)/2}.
\]
Yet the minimum relative gap is $1/2$. For fixed failure probability and
arbitrarily large $k$, this contradicts a lower bound equal to a positive
universal constant times $\operatorname{poly}(k)^{-1}c^{-k}$.
It does not contradict any proved theorem in \cite{chen}, which retains
the condition number, nor the relative convergence bound above.
Subspace geometry can be well behaved when raw monomial coordinates
are badly conditioned.

\section{Limits and remaining review}
This proves a direct bound with Shao's partition-specific gap and
$d-1$ exponent. It does not make the refuted intermediate
$\chi_{\rm mono}\chi_{\rm coef}$ constants bounded, nor remove the need
for cross-block separation. Standard block Krylov, scalar filters,
Gaussian anti-concentration and orthogonal polynomial recurrences are
established tools. The proposed new ingredient is their direct
terminal-resolvent and conditional-row composition, including the
fractional-moment control of growing depth.

No hardware performance claim, new multiplication exponent, unrestricted
quadratic-time matrix-product algorithm, or backward-stability theorem
follows from this manuscript. Numerical monomial Krylov matrices may be
ill-conditioned; the actual algorithm is the usual orthogonalized block
Lanczos method. The theorem is a theoretical convergence statement.
Independent mathematical and primary-literature reviews remain necessary
before a final novelty or significance claim.

\begin{thebibliography}{9}
\bibitem{shao} N. Shao.
\emph{A structural bound for cluster robustness of randomized small-block
Lanczos}. arXiv:2507.10144v2, 30 May 2026.
\url{https://arxiv.org/abs/2507.10144}
\bibitem{stepaniants} G. Stepaniants.
\emph{KE-05: failure of spectrum-uniform probabilistic interpolation
bounds}. OpenProblemsInNLA, September 2026.
\url{https://github.com/ajt60gaibb/OpenProblemsInNLA/blob/main/randomized-and-low-rank-approximation/KE-05/solution.md}
\bibitem{hmt} N. Halko, P.-G. Martinsson, J. A. Tropp.
\emph{Finding structure with randomness: probabilistic algorithms for
constructing approximate matrix decompositions}. SIAM Review 53(2), 2011.
\url{https://arxiv.org/abs/0909.4061}
\bibitem{mmm} R. A. Meyer, C. Musco, C. Musco.
\emph{On the unreasonable effectiveness of single vector Krylov methods
for low-rank approximation}. SODA 2024, pp. 811--845.
\url{https://doi.org/10.1137/1.9781611977912.32}
\bibitem{chen} T. Chen, E. N. Epperly, R. A. Meyer, C. Musco, A. Rao.
\emph{Does block size matter in randomized block Krylov low-rank
approximation?} SODA 2026, pp. 1026--1046.
\url{https://doi.org/10.1137/1.9781611978971.42}
Full version:
\url{https://arxiv.org/abs/2508.06486}
\end{thebibliography}
\end{document}
